%---------------------------------------------------------------------
\chapter{Rigour, Threats to Validity, and Honest Non-Claims}
\label{ch:rigour}

This chapter exists because the campaign's strongest property is not any single
score. It is that the boundary between what was proven, what was measured, and
what remains open is stated explicitly, in the same document as the results, at
the same level of specificity. A report that claims twelve complete suites would
be easier to write and impossible to verify. This one is verifiable.

\section{The evidence rule}

Every number in this report is bound to a surviving artefact in the version
-controlled repository. The rule has four operative clauses:

\begin{enumerate}[leftmargin=1.6em,itemsep=3pt]
  \item \textbf{No pooled denominators.} Original-contract and replacement-scope
  numbers are never combined, and a partial sub-scope is never reported as a
  suite score.
  \item \textbf{No cross-suite averages.} Fourteen suites with different
  semantics do not have a meaningful mean. None is reported.
  \item \textbf{No baseline-gain claims.} Several suites have no measured
  baseline under matched conditions. Where a baseline was not measured, no
  improvement is claimed --- including where the result is complete and strong.
  \item \textbf{Failures stay in the denominator.} A generation failure or an
  artefact loss is retained, even when removing it would improve the headline
  percentage.
\end{enumerate}

Clause~4 is the one that costs the most and therefore proves the most. E1 reports
$2/731$ with 14 generation failures and 4 artefact losses retained. E8 reports a
0.000 category accuracy for CloningScenarios rather than omitting a category that
dragged the profile down. These choices are visible in the numbers themselves.

\section{Honest non-claims}

Stating what is \emph{not} claimed is a load-bearing part of the methodology, not
a hedge. The following are explicit non-claims:

\begin{itemize}[leftmargin=1.4em,itemsep=2pt]
  \item \textbf{No causal capability gain.} No result establishes that training
  improved the model relative to an unmeasured baseline.
  \item \textbf{No security claim.} E10 measures benign task completion. It is not
  a prompt-injection resistance result.
  \item \textbf{No serving parity.} Three serving configurations were used
  (eager, graph at 32k, graph at 65k context). No numerical parity between them is
  claimed.
  \item \textbf{No algorithmic win.} No new GRPO variant is proposed as superior.
  Phase 1's algorithmic levers returned honest nulls and Phase 2 did not reopen
  them.
  \item \textbf{No integrity generalisation.} The Phase-1 P8 method-trace
  classifier is a proof of concept; cross-run integrity detection remains untested.
  \item \textbf{No aggregate grade.} The fourteen lanes do not sum to a
  program-level score, and no such aggregation is offered.
\end{itemize}

\section{Known threats}

The threats below are recorded rather than hidden, in the form a reviewer would
need to assess them.

\begin{itemize}[leftmargin=1.4em,itemsep=3pt]
  \item \textbf{Contamination checks are bounded.} Source-overlap checks cover a
  defined source envelope. They do not cover images, semantic similarity, or
  universal pretraining exposure. Any decontamination statement in this report
  carries that bound.
  \item \textbf{Serving-configuration heterogeneity.} LAB-Bench combined 240
  eager-serving responses, 1504 graph-serving responses at 32{,}768-token context,
  and 223 at 65{,}536-token context. Original answers were retained across
  configurations; the configuration is a source of variance that is disclosed, not
  removed.
  \item \textbf{Evidence-source loss.} The \texttt{.codex-run} deletion removed
  finish-era recovery code. It is confined to build/recovery code; native runners
  survive; sealed requests and hashes pin faithful re-implementation. Reproduction
  is less convenient, not invalid.
  \item \textbf{Small held-out sets in Phase 1.} Held-out $n=8$--$20$ in the
  Phase-1 studies. Magnitudes are noise-limited and are stated as such.
  \item \textbf{Single-seed and single-trajectory Phase-1 results.} P2 is a
  shared-schedule seed-0 diagnostic; P3 is a single-seed sweep; P8 is one
  trajectory per method. Each is labelled exploratory, not confirmatory.
  \item \textbf{A partially recorded sweep artefact.} G8/G16 summaries in the P3
  sweep were lost to a logging thread-safety bug. Re-running under a
  process-per-run launcher yields corrected held-out gains of $+0.083$ (both from
  $0.750\to0.833$, single seed), which \emph{supersede} the affected values --- but
  they remain single-seed point estimates pending multi-seed replication.
  \item \textbf{An invalidated metric from the first campaign batch.} A
  \texttt{zero\_loss\_frac} field was invalid in the first batch and has since been
  fixed. No result in this report depends on it.
  \item \textbf{Judge-instrument degradation.} See below.
\end{itemize}

\section{The judgment-layer incident}

The campaign included a receipt-recording judgment harness used for triage
classification, claim-faithfulness, and value ordering. Its behaviour is reported
because it bears directly on which claims in this report rest on what evidence.

\textbf{Completed before degradation.} The research-value ordering of
Chapter~\ref{ch:p2results} (14 lanes, one ask battery) completed and its receipt
survives.

\textbf{Invalidated by the incident.} Both blocker-classification batteries
returned degenerate distributions --- near-uniform, or a constant class across all
lanes --- after approximately 13:10Z, when both the latest and preview instances
stopped reading state. A controlled yes-probe fell from approximately $1.0$ to
approximately $0.3$. Those receipts are retained as \emph{incident evidence} and
are not used to support any claim. They must be re-run before the judgment layer
can be relied upon again.

\textbf{Consequence for this report.} The per-lane ``first uncleared blocker''
column is grounded in document evidence, \textbf{not} in judgment-layer output.
This is stated explicitly so that a reader knows exactly which claims rest on
which instrument.

\textbf{Design rule banked from the healthy window.} Arithmetic and string
lookups stay in code. The judgment layer once mis-answered a parity probe even
while healthy; it is reserved for classification, faithfulness, and value
ordering, exactly as its documentation intends. A model verdict is never reported
as an empirical result anywhere in this report.

\section{Rigour practices carried from Phase 1}

\begin{itemize}[leftmargin=1.4em,itemsep=2pt]
  \item \textbf{Multi-seed with matched baselines where feasible.} The curriculum
  intervention was re-tested across three seeds at matched token budget; the P1
  scaling check used four seeds.
  \item \textbf{Held-out evaluation with reported $n$.} Every Phase-1 gain is
  reported alongside its held-out $n$ so the reader can see the noise floor.
  \item \textbf{Recompute-from-tensors.} Published diagnostics are recomputable
  from stored raw tensors, not merely read back from a logged summary.
  \item \textbf{Adversarial review plus deterministic verification.} Adversarial
  review of early single-seed conclusions prompted the multi-seed and matched-token
  re-tests. The retained evidence is script-based recomputation: $\zvf$
  $0.72$--$0.77$, all-correct $0.65$--$0.71$, and the stored P8 row-CV output.
\end{itemize}

\section{AI-assistance disclosure}

Language models assisted with adversarial review, code review, and editorial
revision. They were also used as an independent auditor of some claims. These
model assessments are \textbf{not} treated as scientific evidence and are not
required to reproduce any number in this report. Verification in this submission
is script-based, using stored tensors, hashes, and explicit result JSONs. The
candidate remains responsible for the research design, execution choices, source
verification, calculations, and every final claim.
